
\chapter*{Prefacio} 

\justifying La Geometría Proyectiva nació como solución a los problemas de perspectiva que
aparecen al tratar de representar en un plano objetos del espacio. Este tipo de necesidades, surgidas en la pintura y arquitectura, llevaron a la creación de una teoría
matemática que recogiera las bases geométricas subyacentes a los métodos concretos
que poco a poco se fueron elaborando en esos campos del arte.

Las primeras ideas de geometría proyectiva se remontan al Renacimiento motivadas
por la actividad de artistas y arquitectos. Piero della Francesca (1410-1492), Leone
Battista Alberti (1404-1472) y Alberto Durero (1471-1528) reflexionaron sobre nociones
geométricas como las de proyección y sección para entender el problema de la representación de un objeto real tridimensional en un lienzo plano. El primer matemático
que hizo uso de esas ideas fue el francés Girard Desargues (1591-1661), que introdujo los métodos proyectivos en su estudio de las cónicas. Nociones proyectivas como
los puntos del infinito fueron introducidos por Desargues al mismo tiempo que nuevos
métodos de demostración, que le permitieron realizar el mayor avance en el estudio
de las cónicas desde los tiempos de Apolonio. Mientras que los griegos tenían que
idear un método diferente para cada teorema particular, Desargues vio en el enfoque
proyectivo un procedimiento general que le permitía establecer resultados para todas
las cónicas. Desargues influyó fuertemente sobre Blaise Pascal (1623-1662) y Philippe
de la Hire (1640-1718), que dedicaron una actividad considerable al estudio proyectivo
de las cónicas.
 
La Geometría Proyectiva estudia las propiedades de las figuras que se conservan por
proyección. Así, los matemáticos anteriormente mencionados enunciaron nuevos teoremas en los que nociones geométricas clásicas como longitud y ángulo eran sustituidas
por otras de incidencia de las figuras que, a diferencia de las primeras, sí se conservaban por proyección. Progresivamente se fue haciendo claro que los métodos basados en
estas ideas eran mucho más generales que los tradicionales y, a partir de una cuidadosa
distinción entre las propiedades proyectivas y las métricas, a finales del siglo XIX se
configuró una geometría mucho más amplia, la Geometría Proyectiva, que engloba a
todas las geometrías conocidas anteriormente. Esto es, las geometrías afín, euclídea y
no euclídeas se pueden considerar subgeometrías de la proyectiva. Un ejemplo claro
de este hecho es el modelo de geometría hiperbólica que proporciona la Geometría
Proyectiva especialmente adecuado para tener una visión intuitiva de dicha geometría
no euclídea.

Por otro lado, en la Geometría Proyectiva se ha desarrollado la noción de transformación que, a partir de las ideas presentadas por Felix Klein (1849-1925) en su “programa
de Erlangen”, ocupa un lugar central en la construcción de cualquier teoría geométrica.
De esta forma, el concepto de invarianza por transformaciones, que tan importante ha 
sido en la física moderna, ha sido largamente anticipado en la Geometría Proyectiva
al estudiar las propiedades geométricas que permanecen invariantes por cambios de
coordenadas.

Dentro del cuerpo de las matemáticas, la Geometría Proyectiva se puede considerar
como un modelo de teoría que teniendo sus orígenes en el mundo real, llega a un
nivel de elaboración en el que se combinan “la perfección formal, la elegancia en el
razonamiento y la riqueza de las intuiciones”.
Hay que destacar que a principios del siglo XIX, la presentación que se hizo de la
Geometría Proyectiva fue de tipo sintético o axiomático. Una de las figuras más importantes fue Jean-Victor Poncelet (1788-1867), que estudió las figuras homólogas, es decir,
relacionadas entre sí a través de proyecciones y secciones, dando la definición de transformación proyectiva como aquellas aplicaciones que son composición de proyecciones
o perspectividades. Otra figura importante fue Karl G. C. von Staudt (1798-1867),
quien contribuyó de modo significativo a establecer los fundamentos de la Geometría
Proyectiva, haciéndola independiente de las nociones de longitud, distancia y ángulo
y, por tanto, aclarando desde un punto de vista lógico que la Geometría Proyectiva es
anterior a la euclídea. Además, las transformaciones proyectivas que consideró fueron
las aplicaciones biyectivas que transforman 3 puntos alineados en tres puntos alineados,
conocidas también como colineaciones.

La introducción de las ideas algebraicas en geometría proyectiva tuvo lugar también en esas fechas. Sus impulsores fueron Augustus F. Möbious (1790-1868) y Julius Plücker (1801-1868). Ellos fueron los primeros en usar las coordenadas homogéneas para formular algebraicamente muchas ideas proyectivas. A finales del siglo XIX, con
las ideas de la geometría moderna, se consiguió que las demostraciones con cálculos en coordenadas desaparecieran de la geometría proyectiva y se impusieran los razonamientos puramente lógicos a través del álgebra lineal. En estos apuntes, comenzaremos con la presentación algebraica de la Geometría Proyectiva a partir de las rectas que pasan por el origen de un espacio vectorial. Posteriormente veremos la presentación axiomática, quizás más intuitiva desde el punto de vista de una geometría formal, y finalmente comentaremos la equivalencia de las dos presentaciones para el caso del plano proyectivo, construyendo el cuerpo de escalares a partir de propiedades geométricas como el teorema de Desargues y el teorema de Pappus. En esta equivalencia es donde se aprecia la importancia del “Teorema Fundamental de la Geometría Proyectiva”, el cual afirma que las transformaciones de las dos geometrías, algebraica y axiomática, son las mismas.

\hfill Esteban Gómez. 4 de abril de 2003.